\section*{अध्याय \ref{ch:b1:rifts-ridges-margins} --- दरारें, कटक और निष्क्रिय किनारे}

\begin{solution}{exo:b1:rifts-ridges-margins:1}
(क) डाइक मानचित्र पर $\qty{60000}{m}\times\qty{8}{m} = \qty{4.8e5}{m^2}$ घेरता, अतः उसकी औसत ऊँचाई $\qty{2.5e9}{m^3}/\qty{4.8e5}{m^2} =
\qty{5200}{m}$ होती, लगभग \qty{5}{km}: डाइक ने \omterm{def:b1:layered-earth:layers}{भूपर्पटी} के बहुत बड़े भाग को काटा। (ख) $\qty{8}{m}/\qty{0.018}{m.yr^{-1}} = 444$ वर्ष, लगभग चार सदियाँ।
\end{solution}

\begin{solution}{exo:b1:rifts-ridges-margins:2}
(क) \omterm{def:b1:cross-sections:throw}{क्षैतिज विस्थापन} \omterm{def:b1:cross-sections:throw}{विस्थापन} को \omterm{def:b1:geological-maps:attitude}{नति} की स्पर्शज्या से भाग देने पर मिलता है, इस प्रकार:
$3.0/\tan\ang{60} = \qty{1.73}{km}$ और इसी प्रकार $0.8/\tan\ang{70} = \qty{0.29}{km}$;
खंड का विस्तार \qty{2.02}{km} है। (ख) भ्रंशन से पहले खंड $25 - 2.02 = \qty{22.98}{km}$ चौड़ा था, लगभग \qty{23}{km}: वह लगभग \qty{9}{\percent} लंबा हुआ है।
\end{solution}

\begin{solution}{exo:b1:rifts-ridges-margins:3}
अति-धीमे (\qty{20}{mm.yr^{-1}} से कम): गाक्कल कटक (12), मध्य लाल सागर (15) और रेक्यनीस कटक (19), जिनमें अंतिम सीमा के निकट है। धीमे (मान 20 से \qty{55}{mm.yr^{-1}}): पश्चिमी कार्ल्सबर्ग कटक (23) और दक्षिण अटलांटिक (33)। तेज़ (\qty{75}{mm.yr^{-1}} से ऊपर): पूर्वी प्रशांत उभार (143)। अति-धीमे और धीमे कटकों में सामान्य भ्रंशों के बीच एक अक्षीय घाटी होनी चाहिए; पूर्वी प्रशांत उभार में एक अक्षीय उभार।
\end{solution}

\begin{solution}{exo:b1:rifts-ridges-margins:4}
(क) $\qty{3.0}{km^2}\times\qty{7.1}{km} = \qty{21.3}{km^3}$, लगभग
\qty{21}{km^3} प्रति वर्ष। (ख) $\qty{3.028e8}{km^2}/\qty{3.0}{km^2.yr^{-1}} =
\qty{1.01e8}{yr}$, अर्थात् \qty{101}{Myr}। सबसे पुराना तल, लगभग \qty{190}{Myr}, इसलिए पुराना है क्योंकि तल आयु के क्रम में नष्ट नहीं होता: निमज्जन हर आयु का तल खाता है, कुछ पुराना तल वहाँ बचता है जहाँ कोई खाई उस तक नहीं पहुँची, और उत्पादन की दर बदलती रही है।
\end{solution}

\begin{solution}{exo:b1:rifts-ridges-margins:5}
(क) $2500 + 350\times2 = \qty{3200}{m}$, $2500 + 350\times6 = \qty{4600}{m}$,
$2500 + 350\times8 = \qty{5300}{m}$। (ख) $t = \bigl((d - 2500)/350\bigr)^2$:
$(2100/350)^2 = \qty{36}{Myr}$ और $(2800/350)^2 = \qty{64}{Myr}$। (ग) तल 4 और \qty{36}{Myr} के बीच \qty{1400}{m} धँसता है, परंतु 36 और \qty{64}{Myr} के बीच केवल \qty{700}{m}, लगभग समान अवधि में।
\end{solution}

\begin{solution}{exo:b1:rifts-ridges-margins:6}
प्रभावी प्रसार-गुणांक है, वह
$\num{3.2e-5}\times3300/2300 =$ $\qty{4.591e-5}{K^{-1}}$, और
\[ \sqrt{\kappa\times\qty{1}{Myr}} = \sqrt{\num{e-6}\times\num{3.156e13}} = \qty{5618}{m} \]
तब इसके लिए $C =$ $(2/\sqrt\pi)\times\num{4.591e-5}\times1350\times5618 =$
$1.1284\times\num{4.591e-5}\times1350\times5618 = \qty{393}{m.Myr^{-1/2}}$। \qty{250}{K} ठंडे \omterm{def:b1:layered-earth:layers}{प्रावार} के लिए, $C$ $T_m - T_0$ के साथ मापता है:
$393\times1100/1350 = \qty{320}{m.Myr^{-1/2}}$। एक ठंडे महासागर का तल धीमे धँसता है और छोटा स्थिरांक अनुसरण करता है। अध्याय का परिकलित \num{350} ऐसा मामला नहीं है: वह मापदंडों की अनिश्चितता के भीतर है।
\end{solution}

\begin{solution}{exo:b1:rifts-ridges-margins:7}
(क) $\qty{e-6}{cal.cm^{-2}.s^{-1}} = \num{4.184e-6}\times\num{e4} =
\qty{4.184e-2}{W.m^{-2}} = \qty{41.84}{mW.m^{-2}}$, अतः स्थिरांक $11.3\times41.84 = \qty{473}{mW.m^{-2}.Myr^{1/2}}$ है, \qty{3.5}{\percent} नीचे \cref{ch:b1:heat-budget} के 490 से: दोनों फ़िट सहमत हैं। (ख) $q = 473$, $473/\sqrt{10} = 150$ और $473/\sqrt{150} = \qty{38.6}{mW.m^{-2}}$। वर्गमूल नियम लगातार गिरता रहता है, जबकि प्लेट प्रतिरूप एक स्थिरांक पर समतल हो जाता है: \qty{150}{Myr} पर वर्गमूल का मान बहुत कम है, जैसे गहराई-नियम बहुत गहरा चला जाता है।
\end{solution}

\begin{solution}{exo:b1:rifts-ridges-margins:8}
उत्तरी खंड $x = 0$ पर लीजिए और दक्षिणी $x = -40$ पर (\unit{km} में), रूपांतर $y = 0$ पर उनके बीच। (क) रूपांतर के ठीक उत्तर में तल उत्तरी खंड के पश्चिम में है और पश्चिम की ओर चलता है; उसके ठीक दक्षिण में, वह दक्षिणी खंड के पूर्व में है और पूर्व की ओर चलता है। किसी भी ओर से देखने पर दूर का पक्ष बाईं ओर चलता है: सरकाव बाएँ-पार्श्व है। (ख) कटक दाहिने-पार्श्व विस्थापित दिखता है, उसका दक्षिणी भाग पश्चिम में: विपरीत दिशा, जैसा \cref{prop:b1:rifts-ridges-margins:transform} भविष्यवाणी करता है। (ग) भूकंप केवल दो कटक-सिरों के बीच \qty{40}{km} पर होते हैं, जहाँ दोनों पक्ष अलग-अलग प्लेटों के हैं; उससे आगे, भंग-क्षेत्रों के अनुदिश, दोनों पक्ष एक ही प्लेट के हैं और एक साथ चलते हैं।
\end{solution}

\begin{solution}{exo:b1:rifts-ridges-margins:9}
आधी दर \qty{10.6}{mm.yr^{-1}} है, अर्थात् \qty{10.6}{km.Myr^{-1}}। (क) आयु का चरण $350/10.6 = \qty{33.0}{Myr}$ है। (ख) रूपांतर के अंत पर, \qty{33.0}{Myr} पुराना तल ($2500 + 350\sqrt{33.0} = \qty{4511}{m}$) आयु 0 के तल (\qty{2500}{m}) के सामने है: लगभग \qty{2010}{m} का चरण। \qty{100}{km} आगे, आयु $450/10.6 = 42.5$ और $100/10.6 = \qty{9.4}{Myr}$, गहराइयाँ 4781 और \qty{3575}{m}: लगभग \qty{1210}{m} का चरण। \qty{1000}{km} आगे, आयु 127.4 और \qty{94.3}{Myr}, गहराइयाँ 6450 और \qty{5899}{m}: \qty{550}{m}। (ये आयुएँ \qty{80}{Myr} से परे हैं, जहाँ वर्गमूल नियम बहुत गहरा चला जाता है; प्लेट नियम एक छोटा चरण देता।)
\end{solution}

\begin{solution}{exo:b1:rifts-ridges-margins:10}
(क) हर आयु के समान क्षेत्रफल के साथ, माध्य गहराई $d(t)$ का माध्य है, $[0, A]$ पर:
\[
  \frac1A\int_0^A\bigl(d_r + C\sqrt t\bigr)\mathrm dt = d_r + \frac CA\cdot\frac23A^{3/2}
  = d_r + \frac23C\sqrt A.
\]
(ख) $A = \qty{180}{Myr}$ के लिए: $2500 + \frac23\times350\times13.416 =
\qty{5630}{m}$; \qty{120}{Myr} के लिए: $2500 + \frac23\times350\times10.954 =
\qty{5056}{m}$, \qty{574}{m} उथला। (ग) मान लीजिए समुद्र $r$ ऊपर उठता है और लदा तल $s = r\rho_w/(\rho_m - \rho_w)$ धँसता है। पानी, आयतन में अपरिवर्तित, उसी क्षेत्र पर उसी गहराई को भरना चाहिए, अतः $r + s = \qty{574}{m}$,
$r\,\rho_m/(\rho_m - \rho_w) = 574$, और $r = 574\times2300/3300 =
\qty{400}{m}$: केवल प्रसार-दर से सैकड़ों मीटर।
\end{solution}

\begin{solution}{exo:b1:rifts-ridges-margins:11}
(क) $\sqrt{x - a} = \sqrt x\sqrt{1 - a/x} \approx \sqrt x\,(1 - a/2x)$ के लिए
$a \ll x$, अतः
$\Delta d \approx C\sqrt{x/u}\cdot a/(2x) = Ca/(2\sqrt{ux})$।
(ख) $x = \qty{450}{km}$ पर: सटीक \qty{1205}{m}, सन्निकट
$350\times350/(2\sqrt{10.6\times450}) = \qty{887}{m}$; \qty{1350}{km} पर:
\qty{550}{m} और \qty{512}{m}; \qty{3350}{km} पर: \qty{334}{m} और
\qty{325}{m}। सन्निकट और सटीक चरण का अनुपात केवल $r = a/x$ पर निर्भर करता है: $(r/2)/(1 - \sqrt{1 - r})$, जो $r = 0.36$ पर 0.9 है। सन्निकटन \qty{10}{\percent} तक अच्छा है $x = 350/0.36 \approx
\qty{970}{km}$ से आगे, कटक से लगभग तीन विस्थापनों पर।
\end{solution}

\begin{solution}{exo:b1:rifts-ridges-margins:12}
(क) $d_2 - d_1 = 3671.1 - 3606.8 = \qty{64.3}{m}$ पर \qty{10}{Myr};
$4072.8 - 4065.3 = \qty{7.5}{m}$ पर 20; $4415.4 - 4417.0 = \qty{-1.7}{m}$ पर
30; $4707.5 - 4713.6 = \qty{-6.1}{m}$ पर 40। (ख) $d_2 - d_1$ 20 पर धनात्मक और 30 पर ऋणात्मक है, और दोनों नियम सतत हैं: उनके बीच वे बराबर होते हैं। द्विभाजन: 25 पर, $d_2 - d_1 = +0.9$; 27.5 पर, $-0.7$; 26.25 पर, $+0.01$; 26.875 पर, $-0.3$; 26.56 पर, $-0.2$; 26.41 पर, $-0.08$; 26.33 पर, $-0.04$: मूल 26.25 और 26.33 के बीच है, \qty{26.3}{Myr} पर, जहाँ दोनों देते हैं \qty{4294}{m}। (ग) $d_1' = 175/\sqrt{26.3} = \qty{34.1}{m.Myr^{-1}}$ और $d_2' = (3200/62.8)\,e^{-26.3/62.8} = \qty{33.5}{m.Myr^{-1}}$। नियम वहाँ लगभग एक ही मान और एक ही ढलान रखते हैं, अतः जुड़े वक्र में कोई दृश्य मोड़ नहीं है।
\end{solution}

\begin{solution}{pb:b1:rifts-ridges-margins:1}
\textbf{1.} \omterm{def:b1:cross-sections:throw}{क्षैतिज विस्थापन} $= $ \omterm{def:b1:cross-sections:throw}{विस्थापन}$/\tan$(\omterm{def:b1:geological-maps:attitude}{नति}): F1 \qty{346}{m}, F2 \qty{770}{m}, F3 \qty{2014}{m}, F4 \qty{163}{m}, F5 \qty{462}{m}, F6 \qty{1172}{m}। F3, सबसे बड़े \omterm{def:b1:cross-sections:throw}{विस्थापन} और \omterm{def:b1:cross-sections:throw}{क्षैतिज विस्थापन} के साथ, सीमा-भ्रंश है।

\textbf{2.} विस्तार योग है, \qty{4927}{m}, लगभग \qty{4.9}{km}। दरार भ्रंशन से पहले $60 - 4.9 = \qty{55.1}{km}$ चौड़ी थी।

\textbf{3.} \ang{020} पर झुकी दरार का अभिलंब \ang{110} की ओर इंगित करता है; गति, \ang{105} की ओर, उसके साथ \ang{5} बनाती है। लंबवत घटक $6.0\cos\ang{5} = \qty{5.98}{mm.yr^{-1}}$ है।

\textbf{4.} $\qty{4.927}{km}/\qty{5.98}{km.Myr^{-1}} = \qty{0.82}{Myr}$.

\textbf{5.} \qty{3.2}{Myr} में प्लेटें दरार के आर-पार $5.98\times3.2 = \qty{19.1}{km}$ अलग हुईं। भ्रंश उसका $4.93/19.1 = \qty{26}{\percent}$ वहन करते हैं।

\textbf{6.} शेष तीन-चौथाई, \qty{14.2}{km}, परिच्छेद के भ्रंशों में नहीं हैं। उन्हें अक्ष के साथ अंतर्विष्ट डाइकों द्वारा वहन किया जा सकता है, जो दरार को बिना भ्रंशित किए खोलते हैं, जैसे अफ़ार में, या परिच्छेद के बाहर के भ्रंशों द्वारा। \qty{8}{m} चौड़े डाइकों के साथ, $\qty{14200}{m}/\qty{8}{m} \approx 1780$ डाइक चाहिए, औसतन हर $3.2\times10^6/1780 \approx 1800$ वर्ष में एक।

\textbf{7.} C2An से: $(21.0 + 20.2)/2.581 = 41.2/2.581 =
\qty{15.96}{mm.yr^{-1}}$। C3n से: $(41.0 + 40.2)/5.235 = 81.2/5.235 =
\qty{15.51}{mm.yr^{-1}}$।

\textbf{8.} $(41.2 - 2.2)/2.581 = \qty{15.11}{mm.yr^{-1}}$ और
$(81.2 - 2.2)/5.235 = \qty{15.09}{mm.yr^{-1}}$: सुधारने पर, दोनों \qty{0.02}{mm.yr^{-1}} के भीतर सहमत हैं। सुधार छोटे कालखंड के लिए अधिक मायने रखता है, जिसकी दूरी छोटी है।

\textbf{9.} बाहर की ओर \omterm{def:b1:cross-sections:throw}{विस्थापन} एक पार्श्व के दोनों किनारों को एक ही मात्रा से खिसकाता है, इसलिए उनके अंतर में निरस्त हो जाता है। दोनों पार्श्वों पर मिलाकर,
$(81.2 - 41.2)/(5.235 - 2.581) = 40.0/2.654 = \qty{15.07}{mm.yr^{-1}}$।

\textbf{10.} उत्तर-पूर्व, अर्थात् $(41.0 - 21.0)/2.654 = \qty{7.54}{mm.yr^{-1}}$; दक्षिण-पश्चिम: $(40.2 - 20.2)/2.654 = \qty{7.54}{mm.yr^{-1}}$। प्रसार सममित है; पार्श्वों के बीच \qty{0.8}{km} का अंतर, दोनों कालखंडों के लिए समान, एक स्थिर खिसकाव है, जैसे अक्ष को \qty{0.4}{km} केंद्र से हटाकर रखा गया हो।

\textbf{11.} $\cos\Delta = \sin\ang{30.9}\sin\ang{18.0} +
\cos\ang{30.9}\cos\ang{18.0}\cos\ang{16.4} = 0.15870 + 0.78287 = 0.94157$, अतः $\Delta = \ang{19.69}$। $\omega R = 0.403\times(\pi/180)\times6371 =
\qty{44.81}{mm.yr^{-1}}$ के साथ, दर $44.81\sin\ang{19.69} =
\qty{15.10}{mm.yr^{-1}}$ है, प्रश्न 8 और 9 की सुधारी दरें।

\textbf{12.} $\cos\Delta = \sin\ang{30.9}\sin\ang{26.0} +
\cos\ang{30.9}\cos\ang{26.0}\cos\ang{11.4} = 0.22512 + 0.75600 = 0.98112$, अतः $\Delta = \ang{11.15}$ और दर $44.81\sin\ang{11.15} =
\qty{8.7}{mm.yr^{-1}}$।

\textbf{13.} भविष्यवाणियाँ, 15.1 और \qty{8.7}{mm.yr^{-1}}, मापे गए 15 और \qty{9}{mm.yr^{-1}} से मेल खाती हैं। दर $\omega R\sin\Delta$ है, और ध्रुव लाल सागर के उत्तर-पश्चिम में, उत्तर-पूर्वी लीबिया में है: जितना दक्षिण, ध्रुव से उतना दूर, और खुलाव उतना तेज़।

\textbf{14.} $15.09\times5.235 = \qty{79}{km}$ समुद्री तल।

\textbf{15.} $5000/120.6 = \qty{41}{km.Myr^{-1}}$, औसतन \qty{41}{mm.yr^{-1}}।

\textbf{16.} $5000/15.1 = \qty{331}{Myr}$ मध्य दर पर; $5000/8.7 = \qty{575}{Myr}$ उत्तरी दर पर।

\textbf{17.} एक निकटवर्ती ध्रुव के परितः खुलता हुआ, लाल सागर कैंची की जोड़ी की तरह चौड़ा होता है: उसकी चौड़ाई $\sin\Delta$ के अनुपात में बढ़ती है, और वह ध्रुव की ओर संकरा होता एक पच्चर बना रहता है। \qty{8.7}{mm.yr^{-1}} पर, उत्तरी भाग को उन \qty{79}{km} समुद्री तल को बनाने में, जो मध्य भाग के पास अब है, $79/8.7 = \qty{9.1}{Myr}$ लगते हैं।

\textbf{18.} \qty{15}{mm.yr^{-1}} पर लाल सागर की अक्ष एक अति-धीमा कटक है; सामान्य भ्रंशों के बीच एक अक्षीय घाटी अपेक्षित है।

\textbf{19.} $2500 + 350\times1 = \qty{2850}{m}$;
$2500 + 350\sqrt{10} = \qty{3607}{m}$; $2500 + 350\sqrt{40} = \qty{4714}{m}$.

\textbf{20.} $t = (1500/350)^2 = \qty{18.4}{Myr}$.

\textbf{21.} पहले \qty{10}{Myr} में वह $350\sqrt{10} = \qty{1107}{m}$ धँसता है; 70 और \qty{80}{Myr} के बीच, $350(\sqrt{80} - \sqrt{70}) = 350\times(8.944 -
8.367) = \qty{202}{m}$, उसी समय में पाँच गुना कम।

\textbf{22.} $\mathrm dd/\mathrm dt = C/(2\sqrt t) = 175/\sqrt{80} =
\qty{19.6}{m.Myr^{-1}}$, लगभग \qty{0.02}{mm} प्रति वर्ष।

\textbf{23.} $490/\sqrt{80} = \qty{54.8}{mW.m^{-2}}$.

\textbf{24.} वर्गमूल नियम: $2500 + 350\sqrt{80} = 2500 + 350\times8.944 =
\qty{5630}{m}$। प्लेट नियम: $6400 - 3200\,e^{-80/62.8} = 6400 - 3200\times0.2798
= \qty{5505}{m}$। वे \qty{126}{m} भिन्न हैं। \qty{80}{Myr} पर भी दोनों नियम लागू हैं: वर्गमूल नियम लगभग \qty{80}{Myr} तक, प्लेट नियम लगभग \qty{20}{Myr} से; उसके आगे वर्गमूल नियम बहुत गहरा चला जाता है।

\textbf{25.} वर्गमूल नियम \qty{5630}{m} देता है: किलोमीटर के दसवें भाग तक, 80 मिलियन वर्ष के नए महासागर का तल
\textbf{\qty[text-series-to-math=true]{5.6}{km}} की गहराई पर है (प्लेट नियम \qty{5.5}{km} देता है)।
\end{solution}
